\section{Different Initializations for Time Series}
\label{sec:time-series-initialization}

An equation with memory needs initial observations before it can generate its
first new value. These observations form the initial lag history. Changing
that history can change the beginning of a trajectory even when the equation
and future innovations are held fixed. This provides a controlled way to
study how a mechanism responds to different starting conditions and how that
response affects forecasting and explanation evaluation.

\subsection{Initialization Choices}
The current generators use two initialization policies: independent uniform
values within a specified interval, and a joint Gaussian draw from the
stationary law of a reconstructed VAR. The first policy offers direct control
over the starting range. The second preserves the correlations required
between channels and lag positions, following the treatment of initial
observations in \citet[Section~4 and Appendix]{heaps2023stationarity}.

Useful experimental alternatives include a zero history, a constant nonzero
history, manually specified lag values, and values close to the limits of
an allowed domain. These are proposed comparison conditions. The current
high-level trajectory helpers draw their histories automatically; a common
user-supplied-history interface remains to be implemented and validated.
An initialization policy must be checked against the selected generator's
assumptions before it is used.

\subsection{Restrictions and Guarantees for Each Generator}
For the Jacobian method, every initial value must lie in the box used to
verify contraction and domain containment. Changing values within that box
preserves the verified conditions. Enlarging it requires recomputing the
bounds and, potentially, changing the equation's scale; see
Section~\ref{sec:implemented-domain-contraction}.

For Shao--Wu factor saturation, the global conditions permit any fixed finite
history. The influence of that history decreases under the shared-innovation
comparison in \citet[Theorem~5.1 and its proof]{shao2007spectral}. Random
histories need finite second moments for the corresponding second-moment
comparison. For product-output saturation, the amplitude bound must include
the size of the chosen initial history. A larger initial range can increase
the reported bound. Section~\ref{sec:implemented-shao-wu} distinguishes these
two guarantees.

For Heaps' VAR construction, any finite initial history can be propagated by
the stable recurrence. Stationarity from the beginning requires the joint
stationary initialization described in Section~\ref{sec:implemented-heaps}.
An arbitrary fixed history can instead produce a transient: the distribution
changes during the initial portion of the trajectory as the effect of the
starting values fades. The original unconstrained AST accepts finite initial
values through the recurrence evaluator; its subsequent behaviour is assessed
empirically.

\subsection{Controlled Comparisons and Evaluation Windows}
An initialization experiment should keep each sampled mechanism, its
coefficients, and the innovation law fixed while varying its initial history.
Using the same future innovation values across those histories isolates the
effect of initialization. Separate runs with independent future innovations
answer a different question: whether their statistical behaviour becomes
similar despite continued trajectory-level differences.

The returned trajectory arrays retain the initial history as a prefix;
evaluation selects the recursively generated portion that follows it.
The reported window starts with the first recursively generated
observation and includes its transient. Comparisons should report the
starting-history policy, the generated horizon, and early-versus-late
diagnostics explicitly. Repeated trajectories of the same mechanism allow
variation due to initialization and noise to be distinguished from variation
between mechanisms. The explosion, low-amplitude, and tail diagnostics in
Section~\ref{sec:generator-diagnostics} provide a shared basis for these tests.

% Planned study: zero/nonzero fixed histories; independently sampled histories
% at several admissible scales; stationary joint VAR histories; boundary cases.
% Pair identical mechanisms and innovation arrays when isolating initialization.
% Evaluate forgetting curves, ensemble moments, finite-horizon quality metrics,
% forecasting performance, initialization cost, and seed/worker reproducibility.
% Keep theoretical guarantees distinct from observations on finite simulations.
